The resulting formulation is standard. A great treatment is developed in Chapter~7 of \cite{hjelmstad2005fundamentals}.

With general warping:
\begin{equation*}
\hat{\bm{x}}(x,y,z)=\bm{x}(x)+\bm{\theta}(x) \times \bm{r}(y,z)
+\bm{\Phi}(y,z)\bm{\alpha}(x)
\end{equation*}
With only transverse warping:
\begin{equation*}
\hat{\bm{x}}(x,y,z)=\bm{x}(x)+\bm{\theta}(x) \times \bm{r}(y,z)
+\left.\bm{\alpha}(x)\cdot\bm{\varphi}(y,z)\right.\; \mathbf{i}
\end{equation*}

\begin{equation}
\begin{aligned}
\ShearStrain &= \bm{u}^{\prime}+\FrameBasis  \times \bm{\theta} \\ %+ O(\delta^2)
\CurveStrain &= \bm{\theta}^{\prime}% + O(\delta^2)
 %= \theta_x^{\prime} \, \FrameBasis 
 % + \theta_y^{\prime} \FrameBasis_y 
 % + \theta_z^{\prime} \FrameBasis_z
\end{aligned}
\label{eq:basic-shear-strains}
\end{equation}

\begin{equation}
\bm{e}= 
\begin{pmatrix}
\epsilon \\ 
\gamma_y \\
\gamma_z \\
\tau \\
\kappa_z \\ 
\kappa_y \\ 
\end{pmatrix} =
\begin{pmatrix}
u_x^{\prime} \\ 
\left(u_y^{\prime}-\theta_z\right) \\ 
\left(u_z^{\prime}+\theta_y\right) \\
\theta_x^{\prime} \\ 
\theta_y^{\prime} \\
\theta_z^{\prime} \\
\end{pmatrix}
\end{equation}
Note that because the theory is linear one has:
\begin{equation*}
\bm{e}= \begin{bmatrix}
\bm{1} \partial_{\xi} & \FrameBasis^\times \\
 & \bm{1} \partial_{\xi}  \\
\end{bmatrix}\begin{pmatrix}
\bm{u} \\ \bm{\theta}
\end{pmatrix}
\end{equation*}

\subsubsection{Variations}

$$
\bm{A} = \begin{bmatrix}
\bm{1} & \\
& \bm{1}
\end{bmatrix}
$$

\begin{equation}
\EquilibriumDerivative^* = \begin{bmatrix}
\bm{1} \partial_{\xi} & \FrameBasis^\times  \\
& \bm{1} \partial_{\xi}  \\
\end{bmatrix}
\end{equation}
where using $\FrameBasis^\times = \FrameBasis_z\otimes\FrameBasis_y - \FrameBasis_y \otimes\FrameBasis_z$ one has explicitly
\begin{equation*}
\EquilibriumDerivative^* 
= \left[\begin{array}{ccc|ccc}
\partial \\
         & \partial &          &           &  &   -1 \\
         &          & \partial &           &  1 &  \\ \hline
         &          &          & \partial  \\
         &          &          &           & \partial \\
         &          &          &           &          & \partial \\
\end{array}\right]
\end{equation*}

\subsubsection{Equilibrium}

Integration by parts is performed in the standard manner
\begin{equation*}
    \left<\bm{\eta},\EquilibriumDerivative\bm{\sigma}\right> = \left.\bm{\eta}\cdot\bm{\sigma}\right| - \left<\EquilibriumDerivative^*\bm{\eta},\,\bm{\sigma}\right>
\end{equation*}
and one finds the adjoint \(\EquilibriumDerivative\)
\begin{equation}
\EquilibriumDerivative = -\begin{bmatrix}
\bm{1} \partial_{\xi} &  \\
\FrameBasis^\times & \bm{1} \partial_{\xi}  \\
\end{bmatrix}
% = -\left[\begin{array}{ccc|ccc}
% \partial \\
%          & \partial \\
%          &          & \partial \\ \hline
%          &          &          & \partial  \\
%          &          &       -1 &           & \partial \\
%          &        1 &          &           &          & \partial \\
% \end{array}\right]
\end{equation}

\begin{equation*}
-\EquilibriumDerivative\bm{A}_* \bm{s}_{\mathrm{p}} = \left\{\begin{aligned}
& \ShearResult^{\prime} \\
& \CurveResult^{\prime} + \FrameBasis \times\ShearResult
\end{aligned}
\right.
\end{equation*}
Expanding furnishes the familiar expressions
\begin{equation*}
\left.\begin{aligned}
N^{\prime} \\
V_y^{\prime} \\
V_z^{\prime} \\
T^{\prime} \\
M_z^{\prime} + V_y \\
M_y^{\prime} - V_z
\end{aligned}
\right\}
\end{equation*}


\begin{ambox}[Basic Shear Frame]{box:basic-shear}

$$
-\EquilibriumDerivative\bm{A}_* \bm{s}_{\mathrm{p}} = \left\{\begin{aligned}
& \ShearResult^{\prime} \\
& \CurveResult^{\prime} + \FrameBasis \times\ShearResult \\
\end{aligned}
\right.
$$

\paragraph{Kinematics}
$$
\begin{aligned}
\ShearStrain &= \bm{u}^{\prime}+\FrameBasis  \times \bm{\theta} \\ %+ O(\delta^2)
\CurveStrain &= \bm{\theta}^{\prime}% + O(\delta^2)
 %= \theta_x^{\prime} \, \FrameBasis 
 % + \theta_y^{\prime} \FrameBasis_y 
 % + \theta_z^{\prime} \FrameBasis_z
\end{aligned}
$$

\end{ambox}